\chapter{Approach and Testbed}
\label{ca:chap:design}

\section{The training loop}

Figure~\ref{ca:fig:loop} shows one training run. The selector picks an arm $(k, d)$. The
student answers $b = 8$ prompts from domain $d$, sampling $G = 8$ responses to each. Each
response is reduced to a canonical final answer and rewarded by how well it agrees with
the answers provider $k$ gave to the same prompt. The student takes one policy-gradient
step, and the budget is charged $c_k$ for each prompt on which the provider was
consulted. Every 10 rounds the student is evaluated on a held-out selection set and the
result is passed back to the selector. Only prompts and generated text cross the boundary
between the student and a provider.

\begin{figure}[htbp]
\centering
\begin{tikzpicture}[
    box/.style={draw, rounded corners, align=center, minimum width=27mm,
                minimum height=10mm, font=\small},
    arr/.style={-{Stealth}, thick}
  ]
  \node[box] (sel)  at (0, 0)    {Selector\\chooses $(k, d)$};
  \node[box] (prov) at (5.3, 0)    {Provider $M_k$\\(black-box, priced)};
  \node[box] (upd)  at (10.6, 0)   {Agreement reward\\and GRPO step};
  \node[box] (stu)  at (10.6, -2.6) {Student $\pi_\theta$\\(LoRA adapter)};
  \node[box] (eval) at (0, -2.6)  {Selection-set\\accuracy per domain};
  \draw[arr] (sel) -- node[above, font=\scriptsize]{$b$ prompts} node[below, font=\scriptsize]{charge $b\,c_k$} (prov);
  \draw[arr] (prov) -- node[above, font=\scriptsize]{text only} (upd);
  \draw[arr] (upd) -- node[right, font=\scriptsize]{$\theta_{t-1}\!\to\!\theta_t$} (stu);
  \draw[arr] (stu) -- node[below, font=\scriptsize]{every 10 rounds} (eval);
  \draw[arr] (eval) -- node[left, font=\scriptsize]{feedback} (sel);
  \draw[arr, dashed] (stu.north west) -- node[above left, font=\scriptsize,
    pos=0.6]{rollouts} (upd.south west);
\end{tikzpicture}
\caption{The provider-selective training loop.}
\label{ca:fig:loop}
\end{figure}

\section{Tasks, student and providers}

\paragraph{Tasks.} Three domains have answers a program can check: three-number Countdown
puzzles \autocite{patino_gold_2025}, MATH problems at difficulty levels 3 to 5
\autocite{hendrycks_math_2021}, and Python function-writing problems from KodCode
\autocite{xu_kodcode_2025}, kept only if their reference solution passes its tests in our
sandbox. Each domain has a 2,000-prompt training pool, a 500-prompt selection set for the
selector's feedback, and a 500-prompt test set used only for final results. Generated code
runs in a bubblewrap sandbox with no network and a read-only file system.

\paragraph{Student.} The student is Qwen2.5-1.5B-Instruct \autocite{qwen_qwen25_2024},
trained through a rank-16 LoRA adapter \autocite{hu_lora_2022}. Before training it scores 0.208 on
Countdown, 0.474 on MATH and 0.396 on code on the selection set.

\paragraph{Providers.} Nine open models were screened on 200 prompts per domain and five
were kept (Table~\ref{ca:tab:pool}). They are served locally but accessed only through a
text interface, so black-box access is enforced rather than assumed. The pool has a strong
and expensive model, a mid-priced generalist, a code specialist, a cheap model the student
could catch up with, and a cheap model from a second family. A sixth provider is a copy of
Phi-4 at the same price whose answers are corrupted on 30\% of prompts. A selector that
cannot detect the corruption has no reason to avoid it.

\begin{table}[htbp]
\centering
\small
\begin{tabular}{@{}llrccc@{}}
\toprule
 & & & \multicolumn{3}{c}{\textbf{Ceiling $\mu_{k,d}$}} \\
\cmidrule(l){4-6}
\textbf{Provider} & \textbf{Size} & \textbf{Price $c_k$} & \textbf{Countdown} &
  \textbf{MATH} & \textbf{Code} \\
\midrule
Phi-4 \autocite{abdin_phi4_2024}                 & 14.7B & 2.92 & 0.874 & 0.804 & 0.686 \\
Qwen2.5-7B \autocite{qwen_qwen25_2024}           & 7.6B  & 1.00 & 0.810 & 0.750 & 0.580 \\
Qwen2.5-Coder-7B \autocite{hui_qwen25coder_2024} & 7.6B  & 0.98 & 0.620 & 0.678 & 0.612 \\
Qwen2.5-3B \autocite{qwen_qwen25_2024}           & 3.1B  & 0.50 & 0.538 & 0.690 & 0.462 \\
Phi-4-mini \autocite{microsoft_phi4mini_2025}    & 3.8B  & 0.84 & 0.616 & 0.618 & 0.558 \\
Phi-4, degraded                                  & 14.7B & 2.92 & 0.614 & 0.532 & 0.492 \\
\bottomrule
\end{tabular}
\caption{The provider pool. Price is GPU time per response relative to Qwen2.5-7B. The
ceiling is the accuracy of the provider's most common answer over its four cached
responses, on the selection set.}
\label{ca:tab:pool}
\end{table}

\paragraph{Prices and cache.} A provider's price is the GPU time it needs to produce one
response, relative to Qwen2.5-7B, measured during screening. Long-winded providers
therefore cost more per query, as under per-token pricing. Each provider answered every
pool and selection prompt four times, and runs read these cached answers, which makes runs
reproducible and the branch measurements below affordable.

\section{Learning from a black-box provider}

For a prompt $x$, let $n_k(z; x)$ count how many of provider $k$'s four cached answers
equal the canonical answer $z$. A student response with answer $z$ receives
\[
  r_k(z; x) = \frac{n_k(z; x)}{\max_{z'} n_k(z'; x)},
\]
or 0 if it has no valid answer. If the provider answered $\{18, 18, 18, 22\}$, the
student's answers 18, 22 and 7 receive 1, $1/3$ and 0. Rewards become group-relative
advantages as in GRPO \autocite{shao_deepseekmath_grpo_2024}, and the student takes one
gradient step per round. No ground truth enters training.

This reward pulls the student, on each prompt, towards the provider's most common answer.
A student that fully learned from provider $k$ would be correct exactly where that answer
is correct, so the provider's majority-vote accuracy, its ceiling $\mu_{k,d}$, measures its
quality as a teacher better than its single-response accuracy. The ceiling is a per-prompt
argument, not a hard limit: in the pilot the student passed the ceilings of two providers
on Countdown by applying what it learned to prompts the provider got wrong.

\section{Selectors}
\label{ca:sec:selectors}

\paragraph{Headroom with learned efficiency (proposed).} The selector scores each arm by
how far the student is below the provider's ceiling, scaled by a learned estimate of how
much accuracy that headroom buys, and divided by price:
\[
  \hat{v}(k, d) = \max\big(\hat{g}_{k,d},\, 0\big) \cdot
                  \frac{\max\big(\mu_{k,d} - \alpha_d,\ 0.05\big)}{c_k},
\]
where $\alpha_d$ is the student's current selection accuracy on domain $d$. The efficiency
$\hat{g}_{k,d}$ is the realised accuracy gain per round per unit of headroom, pooled per
domain, shrunk towards the domain's estimate, and discounted by 0.95 per block so that it
can track a changing student. A domain where no provider produces progress is learned to
be unproductive for all of them at once. The selector plays a random arm with probability
0.1, so every arm keeps being observed.

\paragraph{Cost-aware Thompson sampling.} A model-free comparison: a discounted Gaussian
posterior over each arm's accuracy gain per round, with the price applied only when
choosing, by playing the arm with the largest sampled gain divided by $c_k$
\autocite{xia_budgeted_ts_2015,raj_dts_2017}.

\paragraph{Reference policies.} Uniform allocation over all 18 arms; each of the five real
providers alone, with domains in rotation; a control that trains on ground-truth
correctness with no provider and no cost for 200 rounds; and supervised fine-tuning (SFT)
on provider responses bought with the same budget. The control shows how much providers
add beyond a verifiable reward; SFT shows what the on-policy update adds.

\section{Measuring marginal value}
\label{ca:sec:branching}

H1 and H2 concern $v(k, d; \theta)$, which is measured directly. The uniform run of each
seed saves checkpoints at 25, 50 and 75\% of its budget, and the untrained student serves
as the 0\% checkpoint. From each checkpoint, a short run trains on a single arm for 20
rounds and is evaluated on the full selection set. Marginal value is the change in mean
selection accuracy divided by the budget the branch spent. With 6 providers, 3 domains,
4 checkpoints and 3 seeds, this gives 216 branches. Branches from the same checkpoint and
seed share the adapter, the prompt order and the sampling seeds, so comparisons between
providers are paired.

\section{Design decisions}
\label{ca:sec:decisions}

The main choices, the alternatives considered and the evidence for each choice are listed
in Table~\ref{ca:tab:decisions}.

\begin{table}[htbp]
\centering
\footnotesize
\renewcommand{\arraystretch}{1.1}
\begin{tabular}{@{}>{\raggedright\arraybackslash}p{2.7cm}
                >{\raggedright\arraybackslash}p{3.9cm}
                >{\raggedright\arraybackslash}p{8.3cm}@{}}
\toprule
\textbf{Decision} & \textbf{Alternatives} & \textbf{Reason for the choice} \\
\midrule
Black-box agreement reward & Logit-based on-policy distillation across tokenizers &
  Cross-tokenizer distillation recovered 0.14 of the teacher gap against 0.66 with a shared
  tokenizer, and its hybrid variant collapsed (Section~\ref{ca:sec:motivating}). \\
Normalised vote share & Most-common-answer indicator; raw vote share; rubric with a verifier &
  Carried as much learning signal as ground truth on MATH and code in screening, kept
  signal on Countdown where providers often give invalid answers, and needs no extra model. \\
Majority-vote ceiling as provider quality & Single-response accuracy; agreement with other providers &
  It is where the reward points the student. Peer agreement was label-free but favoured the
  Qwen family, which made up most of the pool. \\
Countdown (3 numbers), MATH 3--5, KodCode & GSM8K; MBPP; 4-number Countdown &
  GSM8K was saturated (all large providers 94--97\%), MBPP had too few problems, and the
  student solved 0.8\% of 4-number puzzles, leaving almost no gradient. \\
Price = GPU time per response & Parameter count; hosted API prices &
  Measured on our hardware and fixed before training. It also captures verbosity: one
  candidate repeated itself until the length limit. \\
Corrupt whole prompts in the degraded provider & Corrupt each response independently &
  Independent corruption at 30\% lowered the ceiling by only 8\%, because the majority
  usually survives. Per-prompt corruption lowers it by 30\%. \\
Learned efficiency; price applied at choice time & Pilot versions of both selectors &
  In the pilot, the headroom selector stalled on MATH and Thompson sampling spent 55\% on
  the degraded provider (Section~\ref{ca:sec:pilot}). \\
\bottomrule
\end{tabular}
\caption{Main design decisions.}
\label{ca:tab:decisions}
\end{table}
